\documentclass[11pt]{article}
\usepackage[margin=1in]{geometry}
\usepackage{amsmath,amssymb,amsthm}
\usepackage[T1]{fontenc}
\usepackage{lmodern,microtype}
\usepackage[colorlinks=true,urlcolor=blue]{hyperref}
\newtheorem{proposition}{Proposition}
\title{Disproof of Conjecture 00000009650\\\large Zero cross-overlaps in the real symmetric class}
\author{gaochengzhecpu}
\date{October 4, 2026}
\begin{document}
\maketitle
\begin{abstract}
For the source's definition $o_{ij}=\langle u_i,u_j\rangle^2$ using
eigenvectors of distinct eigenvalues of the same matrix, all such overlaps
of a real symmetric matrix are zero. In particular the GOE case named in
the source has a point-mass distribution, zero mean, and zero scaled mean.
This contradicts the claimed density and universal $2/n$ scaling.
\end{abstract}

\section*{The elementary obstruction}
We use the real symmetric class alone, so no interpretation of a square of
a complex inner product is needed. Let $T$ be real symmetric and suppose
$Tu_i=\lambda_i u_i$ and $Tu_j=\lambda_j u_j$, where $\lambda_i\ne\lambda_j$.
Self-adjointness of $T$ gives
\[
 \lambda_i\langle u_i,u_j\rangle
 =\langle Tu_i,u_j\rangle
 =\langle u_i,Tu_j\rangle
 =\lambda_j\langle u_i,u_j\rangle.
\]
Thus $(\lambda_i-\lambda_j)\langle u_i,u_j\rangle=0$ and
\[
                       o_{ij}=\langle u_i,u_j\rangle^2=0.
\]
The argument is exact; it holds for every such matrix and pair. In
particular it applies to every distinct-eigenvalue pair of every GOE
realization, since a GOE matrix is real symmetric.

For a simple-spectrum matrix of size $n\geq2$, set
\[
 P_n=\{(i,j):1\leq i<j\leq n\},\qquad m_n=|P_n|=\binom n2.
\]
The strict upper triangle is the relevant one for the source's explicit
``distinct eigenvalues'' requirement. Its empirical measure and mean are
\[
 \mu_n=\frac1{m_n}\sum_{(i,j)\in P_n}\delta_{o_{ij}}=\delta_0,
 \qquad
 \overline o_n=\frac1{m_n}\sum_{(i,j)\in P_n}o_{ij}=0.
\]
More strongly, for \emph{every} test function $f:\mathbb R\to\mathbb R$,
\[
             \frac1{m_n}\sum_{(i,j)\in P_n}f(o_{ij})=f(0).
\]
If eigenvalues repeat and only the distinct-eigenvalue pairs are retained,
each retained overlap is still zero. GOE has simple spectrum almost surely:
repeated eigenvalues are the zero set of the discriminant polynomial, this
polynomial is nonzero (evaluate it on a diagonal matrix with distinct
entries), and GOE has a density on the space of real symmetric matrices.
Consequently the same empirical law is $\delta_0$ almost surely.

\begin{proposition}
The claimed law in Conjecture 00000009650 fails in the real symmetric class.
\end{proposition}
\begin{proof}
For all $n\geq2$, the mean is $0\ne2/n$. If the mean assertion is intended
as the asymptotic equivalent $\overline o_n\sim2/n$, it also fails, since
$n\overline o_n=0$ for every $n$ and cannot tend to $2$.
Furthermore $\mu_n=\delta_0$ for every $n$, so its limit is $\delta_0$.
A probability measure with a Lebesgue density assigns zero mass to the
singleton $\{0\}$, whereas $\delta_0(\{0\})=1$. Hence this limit cannot
have the asserted density, with exponential tails or otherwise.
\end{proof}

\section*{Scope and interpretation}
The counterargument uses the same-matrix inner products printed in both
language versions, not overlaps between eigenvectors of independent
matrices or non-Hermitian left/right eigenvector condition numbers.
The Bourgade--Dubach paper mentioned in the source concerns nonnormal
Ginibre matrices~\cite{BD}; it does not change the elementary identity above.

If one includes diagonal entries despite the distinct-eigenvalue wording,
normalized eigenvectors give $n$ ones and $n(n-1)/2$ zeros in the inclusive
upper triangle. The empirical measure is then
\[
        \frac{n-1}{n+1}\delta_0+\frac{2}{n+1}\delta_1
        \ \Longrightarrow\ \delta_0.
\]
Its mean is asymptotic to $2/n$, but its limit still has no Lebesgue density.
This optional convention does not rescue the complete conjecture.

\section*{Formalization and reproducibility}
The Lean project proves the orthogonality identity over actual real inner
product spaces. It defines the finite set of strict upper-triangle pairs,
the stated overlaps and empirical test-function averages, and proves the
point-mass identity for every test function. It derives the exact mean
contradiction and the failure of convergence of the scaled mean to $2$ for
\emph{any sequence} of real symmetric simple-spectrum eigensystems.

To exhibit actual objects as well, it constructs, for every size $n\geq2$,
$T_n=\operatorname{diag}(0,1,\ldots,n-1)$ with the standard unit basis.
Lean checks symmetry, distinct eigenvalues, normalization, the eigenvector
equations and the violated mean formula. These diagonal examples are not
presented as random GOE samples: the general orthogonality theorem is what
applies to GOE. The probability-zero discriminant observation, elementary
measure interpretation and inclusive-diagonal remark above are explained
in this manuscript; the formal contradiction already follows from the
pointwise mean theorems. No random-matrix universality theorem is assumed.

The project pins Lean 4.19.0 and Mathlib v4.19.0 at commit
\begin{center}\small\texttt{c44e0c8ee63ca166450922a373c7409c5d26b00b}.\end{center}
Reproduction commands, successful build logs, foundational axiom outputs,
file hashes and a separate adversarial self-review accompany the submission.
The Python script performs exact rational finite checks as a supplement;
the all-dimension proofs and real limits are in Lean.

\begin{thebibliography}{1}
\bibitem{BD} P. Bourgade and G. Dubach,
\emph{The distribution of overlaps between eigenvectors of Ginibre matrices},
Probability Theory and Related Fields 177 (2020), 397--464.
\href{https://arxiv.org/abs/1801.01219}{arXiv:1801.01219}.
\end{thebibliography}
\end{document}
